\subsection{A domain is not a function: InterPro2GO, ARBA and transcription factors}\label{sec:dd-domain}
\textbf{Claim.} Signature-based mappings attach a term valid for some family members to all;
association rules (ARBA) OR together thousands of condition sets under generic terms.
\textbf{Corpus.} InterPro2GO is the most \act{MODIFY}-heavy pipeline (Table~\ref{tab:iea}): the
domain is right about the kind of function and wrong about its specificity.
Fig.~\ref{fig:ddpanels}d breaks InterPro2GO, ARBA and TreeGrafter down by aspect: InterPro2GO MF rows
are modified about a third of the time, and TreeGrafter BP has the highest removal rate.
\textbf{Projects.} \proj{INTERPRO.md}: 3,652 judged InterPro2GO rows on 1,706 genes; 36 proposed mapping
edits over 15 entries (e.g.\ ATP binding and protein phosphorylation on IPR000719).
\proj{ARBA\_REVIEW.md}: 6,051 ARBA rows; of those not kept, 62\% true but uninformative (construction
failures) vs 9\% false (consequent failures). \proj{TRANSCRIPTION\_FACTORS.md}: a DNA-binding domain
does not make a transcription factor; GOA agrees with the GREEKC curated dbTF set on 95.6\% of
1,447 proteins, and the GOA-only cases (netrins, HDAC4, RPS3, RFXAP/RFXANK) are errors.
\proj{NCBIFam.md}, \proj{TREEGRAFTER.md}.
\textbf{Upstream.} Curated rules (\texttt{rules/curated/}) that couple sequence features with taxon or
module context. \todo{status of InterPro proposals.}
